\chapter*{考研三卷种阅读路线}
\addcontentsline{toc}{chapter}{考研三卷种阅读路线}
数学一、二、三的线性代数均围绕行列式、矩阵、向量、方程组、特征值与相似、二次型组织训练。优先完成本书第1—6章中对应大纲的概念与例题，再从附录\ref{app:exam_core}回查薄弱点。各卷线性代数在卷面中均约占20\%，不能据此保证某一具体题型每年出现。

数学一还要求 $n$ 维向量空间、子空间、基、维数、坐标与基变换等相关内容，可从第3章和第7章的相应小节回查。这不意味着抽象线性变换理论、商空间、对偶空间都成为数一统考要求。数二、数三仍要掌握方程组的解结构与基础解系，不能因为不单列向量空间专题就跳过齐次解的线性组合。

\section*{先练熟三个已有入口}
\begin{enumerate}
\item \textbf{消元与自由变量：}第\ref{ch:systems}章例\ref{ex:4-5}、\ref{ex:4-6}。先把增广矩阵化为阶梯形，确认相容性，再选自由变量逐一赋单位值。这样生成的是一组线性无关的解向量，一般并不正交；正交化是另一步操作。
\item \textbf{$AB=0$：}回查第\ref{ch:systems}章的乘积与解结构例题。逐列理解 $B$ 的列是 $Ax=0$ 的解，核对乘积中间的公共维数，再写秩关系。此基础结论可以作为统考训练，不必先学完第8章的高阶秩不等式证明。
\item \textbf{相似判定：}第\ref{ch:eigen}章。先以特征多项式、迹、行列式、秩等排除不可能；若两矩阵都可对角化，再核对带重数的特征值。一般矩阵的相似与对角化也在基础范围，不能把考试范围缩成只考实对称矩阵。必要不变量相同不构成一般的充分条件。
\end{enumerate}

\section*{提高内容怎样跳读}
第7章的商空间与对偶、抽象映射结构，第8章的 Jordan 链计数、SVD、伪逆、矩阵范数与条件数等，保留为提高阅读入口，不列入本轮统考主线。不要按“第8章全不看”删除与基础秩、矩阵幂和解结构有关的已有简单题；是否采用提高方法，与题目的结论是否可由基础方法完成，是两个问题。

伴随矩阵的谱与迹速算见第\ref{r2:adjugate}节，非对称矩阵误用正定判据的反例见第\ref{r2:positive}节。大纲版本为教育部教育考试院《2027年全国硕士研究生招生考试数学考试大纲》，核读原书第5、16—19页及数二、数三对应线代范围。
官方版本入口：\href{https://xuanshu.hep.com.cn/front/book/findBookDetails?bookId=6aa1910be119ac97297a608f}{高等教育出版社：2027数学考试大纲}。
